% ===== §1 =====
\section{Introduction: The First Fissure and the Problem of \emph{Wen} and \emph{Hua}}
\label{sec:intro}

\noindent
A previous paper of this series closed not in synthesis but in a set of directed fissures, and the first of them pointed here. That paper concerned the aesthetics of the quotidian and the education of perception, and it argued that between two people turned toward a shared world a capacity for perceiving beauty may be cultivated that neither carries alone \citep{paperxxiii}. It ended by observing that such a shared capacity, once it takes on a durable and transmissible form, has already crossed a boundary the aesthetic alone cannot name. When a gesture repeated between two people hardens into a rite, when a mistaken word is kept on purpose and becomes a name, when a way of seeing is sedimented into a way of living that can be handed on, something has emerged that is no longer a perception and not yet an institution. That something is culture, and this paper takes its emergence in intimate relation for its subject.

The word chosen for the subject is not innocent, and the paper begins by refusing to let it settle too quickly. In English \emph{culture} carries an agricultural inheritance, from \emph{colere}, to till, to tend, to cultivate; the metaphor is of a ground worked and a growth husbanded, and it leans, for all its life, toward the image of a thing cultivated and then possessed, a heritage held. The Chinese term for culture, \emph{wenhua} (\zh{文化}), carries a different charge, and the difference is instructive rather than decorative. It is built of two characters. \emph{Wen} (\zh{文}) is pattern, marking, the woven grain of a thing, the sign; \emph{hua} (\zh{化}) is transformation, generation, the bringing-to-live and bringing-to-form. \emph{Hua} is verbal before it is nominal. To \emph{hua} is to transform, to generate, to make live and live again, and the classical locus---``observe the patterns of the human, and by them bring the world to form''---reads culture not as a deposit but as an ongoing act of formation \citep{legge1963yijing}. The tension between these two inheritances, the tended possession and the ongoing transformation, is not a matter of etymological curiosity. It is the whole problem of the paper in miniature, and it already announces the thesis: that culture is misconceived when it is taken as a finished and transmissible object, and rightly conceived when it is taken as a generation that does not close.

\subsection*{The question, and the special site at which it is asked}

The general question is what culture is, ontologically and epistemologically: in what manner it exists, how it comes to be, and how it can be known. That question has a long literature, and the first part of this paper works through it, not to survey it but to locate a position within it that the existing accounts have left open. But the general question is not, for this series, the final one. This is a philosophy of intimacy, and its concern with culture is a concern with culture at its smallest and most exposed, in the relation between two people. The intimate dyad is the minimal unit of cultural emergence, the smallest system in which a culture can arise at all, and it is therefore the site at which the mechanism of emergence is most nakedly visible. What in a nation or an epoch is buried under the weight of institutions, laws, and the dead accumulation of centuries stands, in a couple, close enough to be watched as it happens. One can almost see the moment a private word is coined and accepted, the moment a shared silence acquires a meaning; and one can also see, with a clarity no macro-culture affords, what it costs when such a culture collapses. The particular is not here an application of the general arrived at late. It is the place where the general is legible.

This double structure organizes the paper. The general and the special are related as universality and the particularity that both instances and exceeds it. The first part develops a general ontology and epistemology of cultural emergence, and it is deliberately general: its claims about how culture comes to be, its dialectic of the material and the ideal, its account of culture across the registers of experience, are meant to hold of culture as such. The second part, which is the heart of the paper and the reason the series exists, descends into the intimate dyad and asks what the general framework becomes when the cultural unit is not a people but a pair, and it finds there mechanisms that belong to the dyad alone and appear nowhere else: the vacancy, between two people, of any external authority to fix meaning; the mutual implication of two structures of desire; a density of shared living that lets culture sediment at a visible speed; a fragility that can return a whole shared world to ruin. The third part turns from what culture is and how it is lived to how it is produced and how it is captured, tracing the material conditions under which intimate culture is generated and the means by which capital, having exhausted older frontiers, reaches into the relation itself to draw value from the very bond.

\subsection*{Three lines run through the whole}

The argument is carried by three lines that are not confined to separate chapters but run through the whole, and it is worth naming them at the outset so that their recurrence is read as design rather than repetition. The first is the line of \emph{historical materialism}: that culture, however much it exceeds its conditions, is produced under material conditions and cannot be understood apart from the forces and relations of its production. This line insists throughout that the emergence of culture is constrained, that its forms are historical rather than eternal, and that the manner of its production and of its capture is a question of political economy and not only of meaning. The second is the line of \emph{the Hegelian dialectic}: that culture is spirit exteriorized, going out of itself into a visible form, coming to know itself in that form, and returning to itself altered, so that exteriorization is not loss but the necessary path of self-knowledge, and so that the movement of culture, including its estrangements, has the shape of a determinate negation rather than a mere succession. The third is the line of \emph{psychoanalysis}: that culture is bound to drive, that it is the coding, the sublimation, the sedimented shared fantasy through which what cannot be directly said is given a transmissible form, and that in intimate relation, above all, culture is the sediment of two structures of desire woven into one. These three lines converge, in the meta-theoretical statement that opens the paper's second movement, on a single framework, a variant of the Lacanian theory of the three registers reworked so that the Real is granted a relational form, so that culture is shown to span all three registers, and so that the registers are set into dialectical and materially driven motion rather than left as a static topology. That framework, named in its place the \emph{three-register dialectic of the relational Real}, is the theoretical spine on which the whole rests, and the introduction names it here only as a promise to be redeemed.

\subsection*{What the paper contributes, with its limits showing}

The contribution can be stated plainly and with its limits in view. It is, first, an ontology of cultural emergence that locates culture neither in a transmissible object nor in a symbolic system nor in an encoding of power alone, but in the dialectical exteriorization, across three registers and under material constraint, of the interwoven generative history of relational subjects. It is, second, a phenomenology of intimate culture that takes the dyad seriously as a special site and describes the mechanisms proper to it, and this is where the paper means to earn its place in the series, for these mechanisms have not, to the author's knowledge, been gathered and named. It is, third, an ethics and a political economy of intimate culture that shows how the value such culture generates is produced, how it is captured, and on what its just generation depends, joining this to the theory of generative justice the series has developed in exchange with Eglash \citep{eglash2016}. The paper does not offer a completed system, and it does not close its own central concept, which would contradict in its form the generativity the concept names. It offers a framework, a set of mechanisms, and a criterion, and it holds them open, in the manner its subject demands, toward what a later work must take up.
